\chapter{Machine models}
\label{chap:models}

This chapter fixes the machine model that every complexity class in the library
is defined from, and the tools for reasoning about it. The base model is the
multi-tape Turing machine of Arora and Barak over the alphabet
$\{0, 1, \sqcup, \triangleright\}$, with a read-only input tape, $k$ work tapes,
and an output tape; deterministic, nondeterministic, and probabilistic machines
share one configuration type and differ only in their transition functions and
acceptance conditions. Formalized so far: the model with its time and output
semantics; machine combinators with Hoare-style rules for sequential,
conditional, looping, and complementing composition, and deciders for
complements and unions; the quadratic reduction to one work tape; a fixed
six-work-tape universal machine proved universal and polynomially efficiently
universal; a weak deterministic time hierarchy; and a logarithmic-cost random
access machine whose polynomial-time class equals $\cc{P}$. The main open
directions are undecidability of the halting problem (now within reach, since
the universal machine is proved to reflect as well as preserve halting), a time
hierarchy under the textbook time-constructibility hypothesis, a
$T \log T$ simulation of Turing machines by random access machines, and
space-level and variant-model robustness for random access machines.

\section{Turing machines}

\begin{definition}[Alphabet and tapes]
  \label{def:models-tape}
  \lean{Complexity.Γ, Complexity.Γw, Complexity.Dir3, Complexity.Tape, Complexity.Tape.HasOutput}
  \leanok
  The read alphabet is $\Gamma = \{0, 1, \sqcup, \triangleright\}$ and the write
  alphabet is $\Gamma_w = \{0, 1, \sqcup\}$, so no transition can write the
  start symbol $\triangleright$. A head moves left, right, or stays. A tape is a
  head position $h \in \mathbb{N}$ together with contents
  $\mathbb{N} \to \Gamma$; the type allows any symbol in any cell. Initial
  tapes have $\triangleright$ in cell $0$ and nowhere else, writing at cell $0$
  is a no-op, and moving left at cell $0$ stays there, so every tape reachable
  from an initial configuration keeps this marker invariant (statements about
  arbitrary tapes assume it explicitly). A tape \emph{has output}
  $y \in \{0,1\}^*$ when cells $1, \dots, |y|$ hold the bits of $y$ and cell
  $|y| + 1$ is blank.
\end{definition}

\begin{definition}[Language]
  \label{def:language}
  \lean{Complexity.Language}
  \leanok
  A language is a set of finite binary strings, $L \subseteq \{0,1\}^*$.
\end{definition}

\begin{definition}[Configuration]
  \label{def:models-cfg}
  \lean{Complexity.Cfg, Complexity.Cfg.init}
  \leanok
  \uses{def:models-tape}
  A configuration of a machine with state set $Q$ and $k$ work tapes consists of
  a state, an input tape, $k$ work tapes, and an output tape, each a named
  field. The initial configuration on input $x$ is in the start state, with $x$
  written from cell $1$ of the input tape, every other tape empty, and every
  head at cell $0$.
\end{definition}

\begin{definition}[Deterministic Turing machine]
  \label{def:tm}
  \lean{Complexity.TM, Complexity.TM.step, Complexity.TM.reachesIn, Complexity.TM.reaches}
  \leanok
  \uses{def:models-cfg}
  A deterministic Turing machine with $k$ work tapes consists of a finite state
  type $Q$ with decidable equality, a start state and a halt state, and a
  transition function
  \[
    \delta : Q \times \Gamma \times \Gamma^k \times \Gamma \to
      Q \times \Gamma_w^k \times \Gamma_w \times \{\mathrm{L}, \mathrm{R}, \mathrm{S}\}^{k+2}
  \]
  that must move right every head reading $\triangleright$. One step applies
  $\delta$ to the current state and the symbols under the input head, the $k$
  work heads, and the output head, writes on the work and output tapes, and
  moves all heads; the halt state has no successor. The relation
  $c \to^t c'$ holds when $c'$ is reached from $c$ in exactly $t$ steps, and
  $c \to^* c'$ is its reflexive--transitive closure.
\end{definition}

\begin{definition}[Nondeterministic and probabilistic Turing machine]
  \label{def:ntm}
  \lean{Complexity.NTM, Complexity.NTM.trace, Complexity.TM.toNTM}
  \leanok
  \uses{def:tm}
  A nondeterministic Turing machine has the same data as a deterministic one
  except that it has two transition functions $\delta_0, \delta_1$, each
  subject to the same rule at $\triangleright$. Its
  execution is defined only through the trace: given a choice sequence
  $r \in \{0,1\}^T$, run $T$ steps, using $\delta_{r_i}$ at step $i$, and stay
  put once the halt state is reached. The same structure serves as a
  probabilistic machine, with $r$ read as random bits. A deterministic machine
  embeds by taking $\delta_0 = \delta_1 = \delta$.
\end{definition}

\begin{definition}[Deciding in time]
  \label{def:tm-decides}
  \lean{Complexity.TM.DecidesInTime}
  \leanok
  \uses{def:tm, def:language}
  A deterministic machine $M$ decides $L$ in time $T : \mathbb{N} \to \mathbb{N}$
  if on every input $x$ it reaches the halt state within $T(|x|)$ steps, with
  output cell $1$ equal to $1$ if $x \in L$ and equal to $0$ if $x \notin L$
  (Arora--Barak Definition~1.2).
\end{definition}

\begin{definition}[Computing a function in time]
  \label{def:tm-computes}
  \lean{Complexity.TM.ComputesInTime, Complexity.TM.Computes}
  \leanok
  \uses{def:tm}
  A deterministic machine $M$ computes $f : \{0,1\}^* \to \{0,1\}^*$ in time $T$
  if on every input $x$ it halts within $T(|x|)$ steps and its output tape has
  output $f(x)$. It computes $f$ if it does so for some $T$.
\end{definition}

\begin{definition}[Nondeterministic deciding]
  \label{def:ntm-decides}
  \lean{Complexity.NTM.DecidesInTime, Complexity.NTM.AcceptsInTime, Complexity.NTM.AllPathsHaltIn}
  \leanok
  \uses{def:ntm, def:language}
  A nondeterministic machine $N$ accepts $x$ in time $T$ if some choice sequence
  of length $T$ leads to the halt state with output cell $1$ equal to $1$. It
  decides $L$ in time $T$ if every trace on every input $x$ has halted after
  $T(|x|)$ steps, and $x \in L$ exactly when $N$ accepts $x$ in time $T(|x|)$
  (Arora--Barak Definition~2.1). Rejection is the absence of an accepting path;
  it does not require output $0$.
\end{definition}

\begin{definition}[Acceptance probability]
  \label{def:models-accept-prob}
  \lean{Complexity.NTM.acceptCount, Complexity.NTM.acceptProb, Complexity.NTM.outputProb}
  \leanok
  \uses{def:ntm}
  For a machine $N$, input $x$, and clock $T$, the acceptance probability is the
  rational number
  \[
    \Pr[N \text{ accepts } x \text{ in } T] =
      \frac{|\{ r \in \{0,1\}^T : \text{the trace halts with output cell } 1 = 1\}|}{2^T},
  \]
  and the probability of output $y$ is defined the same way with the output
  condition in place of acceptance. Both are meaningful when every trace has
  halted by time $T$.
\end{definition}

\begin{lemma}[Deterministic machines as nondeterministic ones]
  \label{lem:models-to-ntm}
  \lean{Complexity.TM.toNTM_decidesInTime, Complexity.TM.toNTM_accepts_iff}
  \leanok
  \uses{def:tm-decides, def:ntm-decides}
  If a deterministic machine decides $L$ in time $T$, its nondeterministic
  embedding decides $L$ in time $T$. Without a time bound, a deterministic
  machine and its embedding accept the same strings, where acceptance means
  reaching the halt state (after some number of steps, along some choice
  sequence) with output cell $1$ equal to $1$.
\end{lemma}
\begin{proof}
  \leanok
  Every trace of the embedding is the deterministic run, frozen at the halt
  state; for acceptance without a bound, the trace of the embedding reaches
  exactly the configurations the deterministic run reaches.
\end{proof}

\section{Output semantics}

\begin{definition}[Halting and producing an output]
  \label{def:produces}
  \lean{Complexity.TM.Produces, Complexity.TM.ProducesInTime, Complexity.TM.Halts, Complexity.TM.HaltsInTime, Complexity.TM.runCfg}
  \leanok
  \uses{def:tm}
  A deterministic machine $M$ halts on a program $p$ (within time $t$) if the
  run from the initial configuration on $p$ reaches the halt state (within $t$
  steps). It produces $y$ on $p$ (within time $t$) if moreover the output tape of
  the halted configuration has output $y$. The bounded evaluator
  $\mathrm{run}(c, t)$ is the configuration after $t$ steps, standing still once
  halted.
\end{definition}

\begin{lemma}[Basic facts about production]
  \label{lem:models-produces-basic}
  \lean{Complexity.TM.ProducesInTime.iff_runCfg, Complexity.TM.instDecidableProducesInTime, Complexity.TM.ProducesInTime.mono, Complexity.TM.produces_iff_exists_producesInTime, Complexity.TM.ProducesInTime.take_time, Complexity.TM.ProducesInTime.output_unique, Complexity.TM.Produces.output_unique}
  \leanok
  \uses{def:produces}
  Let $M$ be a deterministic machine.
  \begin{enumerate}
    \item $M$ produces $y$ on $p$ within time $t$ if and only if
      $\mathrm{run}(\mathrm{init}(p), t)$ is halted with output $y$; in
      particular bounded production is decidable.
    \item Bounded production is monotone in $t$, and $M$ produces $y$ on $p$
      exactly when it does so within some finite time.
    \item If $M$ produces $y$ on $p$ within time $t$, it produces $y$ on the
      first $t$ bits of $p$ within time $t$.
    \item If $M$ produces both $y$ and $y'$ on $p$, then $y = y'$; this holds
      also for bounded production under two different clocks.
  \end{enumerate}
\end{lemma}
\begin{proof}
  \leanok
  Relate the relational run to the bounded evaluator; a run of $t$ steps
  cannot move the input head past cell $t$.
\end{proof}

\begin{lemma}[Computing is pointwise production]
  \label{lem:models-computes-iff-produces}
  \lean{Complexity.TM.computesInTime_iff_forall_producesInTime, Complexity.TM.computes_iff_forall_produces}
  \leanok
  \uses{def:tm-computes, def:produces}
  $M$ computes $f$ in time $T$ if and only if $M$ produces $f(x)$ on $x$ within
  time $T(|x|)$ for every $x$. $M$ computes $f$ if and only if $M$ produces
  $f(x)$ on every $x$.
\end{lemma}
\begin{proof}
  \leanok
  The first equivalence unfolds the definitions. For the second, take the maximum
  halting time over the finitely many inputs of each length.
\end{proof}

\section{Combinators and Hoare-style specifications}

\begin{definition}[Time-bounded Hoare triple]
  \label{def:models-hoare}
  \lean{Complexity.TM.HoareTime, Complexity.TM.Hoare, Complexity.TapePred}
  \leanok
  \uses{def:tm}
  For predicates $\mathit{pre}, \mathit{post}$ on (input tape, work tapes,
  output tape) and a bound $b \in \mathbb{N}$, the triple
  $\{\mathit{pre}\}\, M\, \{\mathit{post}\}_b$ holds if $M$, started in its
  start state on any tapes satisfying $\mathit{pre}$, halts within $b$ steps on
  tapes satisfying $\mathit{post}$. The unbounded triple drops the step bound.
\end{definition}

\begin{definition}[Machine combinators]
  \label{def:models-combinators}
  \lean{Complexity.TM.seqTM, Complexity.TM.ifTM, Complexity.TM.loopTM, Complexity.TM.complementTM, Complexity.TM.unionTM}
  \leanok
  \uses{def:tm}
  On machines with the same $k$ work tapes: $\mathrm{seq}(M_1, M_2)$ runs $M_1$
  to completion and then $M_2$ on the same tapes; $\mathrm{if}(M_t, M_1, M_0)$
  runs the test $M_t$, rewinds the output head, and runs $M_1$ if output cell
  $1$ holds $1$ and $M_0$ otherwise; $\mathrm{loop}(M_b, M_t)$ alternates the
  body and the test, rewinding the output head after each test, and halts once
  the test leaves $1$ in output cell $1$; $\mathrm{complement}(M)$ runs $M$,
  rewinds the output head, and exchanges $0$ and $1$ in output cell $1$ (a
  blank stays blank). In $\mathrm{seq}$ (from $M_1$ to $M_2$), in
  $\mathrm{loop}$ (from the body to the test), and in $\mathrm{if}$ (from the
  branch to the final halt state), the handover is a single step in which
  every head reading $\triangleright$ moves one cell right, the other heads
  stay, and no tape cell changes except that a $\triangleright$ outside cell
  $0$ would be overwritten by a blank. After the test in $\mathrm{if}$ and
  $\mathrm{loop}$ the handover takes several steps: one such step; then steps
  that move the output head left one cell at a time until it reads
  $\triangleright$, and then one cell right; and a last step that reads the
  output cell under the head (cell $1$ when $\triangleright$ appears on the
  output tape only at cell $0$) and enters the chosen branch, the body, or the
  halt state. In each of these steps every input or work head reading
  $\triangleright$ moves one cell right and the others stay, and cells change
  only as in the single-step handover. The machine
  $\mathrm{union}(M_1, M_2)$, with $M_i$ on $k_i$ work tapes, has
  $k_1 + 1 + k_2$ work tapes: it runs $M_1$ with its output redirected to work
  tape $k_1$, halts with $1$ in output cell $1$ if cell $1$ of that work tape
  holds $1$, and otherwise rewinds the input and runs $M_2$ on the last $k_2$
  work tapes and the real output tape.
\end{definition}

\begin{lemma}[Deciders for complements and unions]
  \label{lem:models-complement-union}
  \lean{Complexity.TM.complementTM_decidesInTime, Complexity.TM.unionTM_decidesInTime, Complexity.TM.complementTM_hoareTime}
  \leanok
  \uses{def:models-combinators, def:tm-decides, def:models-hoare}
  \begin{enumerate}
    \item If $M$ decides $L$ in time $T$, then $\mathrm{complement}(M)$
      decides $\{0,1\}^* \setminus L$ in time $2T(n) + 4$.
    \item If $M_1$ decides $L_1$ in time $T_1$ and $M_2$ decides $L_2$ in time
      $T_2$, then $\mathrm{union}(M_1, M_2)$ decides $L_1 \cup L_2$ in time
      $10\,T_1(n) + T_2(n)$.
    \item If $\{\mathit{pre}\}\, M\, \{\mathit{post}\}_b$, where
      $\mathit{post}$ guarantees that the output tape holds $\triangleright$
      at cell $0$ and nowhere else, that the output head is at a cell at most
      $p$, and that output cell $1$ holds a symbol satisfying $\varphi$, then
      $\{\mathit{pre}\}\, \mathrm{complement}(M)\, \{\mathit{post}'\}_{b + p + 4}$,
      where $\mathit{post}'$ says that output cell $1$ holds the flip of some
      symbol satisfying $\varphi$ ($0$ and $1$ exchanged, anything else made
      blank).
  \end{enumerate}
\end{lemma}
\begin{proof}
  \leanok
  For the complement, simulate $M$, then rewind the output head to
  $\triangleright$, step to cell $1$, and write the flipped symbol. For the
  union, simulate $M_1$, rewind its redirected output and read its verdict, and
  either accept at once or rewind the input and simulate $M_2$; the transition
  costs at most $2T_1(n) + 7$ steps, which is absorbed into $9\,T_1(n)$ because
  a decider cannot halt in zero steps (its start state would then be its halt
  state, leaving a blank verdict cell).
\end{proof}

\begin{theorem}[Sequential composition rule]
  \label{thm:models-seq-hoare}
  \lean{Complexity.TM.seqTM_hoareTime}
  \leanok
  \uses{def:models-hoare, def:models-combinators}
  If $\{\mathit{pre}\}\, M_1\, \{\mathit{mid}\}_{b_1}$, the phase-boundary
  normalization maps tapes satisfying $\mathit{mid}$ to tapes satisfying
  $\mathit{mid}'$, and $\{\mathit{mid}'\}\, M_2\, \{\mathit{post}\}_{b_2}$, then
  $\{\mathit{pre}\}\, \mathrm{seq}(M_1, M_2)\, \{\mathit{post}\}_{b_1 + 1 + b_2}$.
\end{theorem}
\begin{proof}
  \leanok
  Simulate $M_1$ in the first phase, take one transition step, and simulate
  $M_2$ in the second phase.
\end{proof}

\begin{theorem}[Conditional rule]
  \label{thm:models-if-hoare}
  \lean{Complexity.TM.ifTM_hoareTime}
  \leanok
  \uses{def:models-hoare, def:models-combinators}
  Suppose $\{\mathit{pre}\}\, M_t\, \{\mathit{mid}\}_{b_t}$, where $\mathit{mid}$
  guarantees that every tape holds $\triangleright$ at cell $0$ and nowhere
  else and that the output head is at a cell at most $p$; that $\mathit{mid}$
  with output cell $1$ equal to $1$ (resp.\ not $1$) implies the precondition
  of $M_1$ (resp.\ $M_0$) on the tapes obtained by normalizing the input and
  work tapes and placing the output head at cell $1$; that
  $\{\mathit{mid}_i\}\, M_i\, \{\mathit{post}_i\}_{b_i}$ for $i = 0, 1$; and that
  both postconditions imply $\mathit{post}$ after normalization. Then
  $\{\mathit{pre}\}\, \mathrm{if}(M_t, M_1, M_0)\, \{\mathit{post}\}_b$ with
  $b = b_t + p + \max(b_1, b_0) + 5$.
\end{theorem}
\begin{proof}
  \leanok
  Simulate the test, rewind the output head to read the verdict, then
  simulate the selected branch.
\end{proof}

\begin{theorem}[Loop rule]
  \label{thm:models-loop-hoare}
  \lean{Complexity.TM.loopTM_hoareTime}
  \leanok
  \uses{def:models-hoare, def:models-combinators}
  Let $\mathit{inv}$ be an invariant and $v$ a variant with $v \le K$ under
  $\mathit{inv}$. Suppose every iteration of $\mathrm{loop}(M_b, M_t)$ started
  from $\mathit{inv}$ either halts within $b$ steps satisfying $\mathit{post}$,
  or returns to the loop start within $b$ steps satisfying $\mathit{inv}$ with
  $v$ strictly smaller. Then
  $\{\mathit{inv}\}\, \mathrm{loop}(M_b, M_t)\, \{\mathit{post}\}_{(K+1) b}$.
\end{theorem}
\begin{proof}
  \leanok
  Induction on the variant.
\end{proof}

\section{Reduction to one work tape}

\begin{theorem}[Single-tape reduction, nondeterministic]
  \label{thm:models-ntm-single-tape}
  \lean{Complexity.NTM.singleTapeSim_decidesInTime, Complexity.NTM.exists_singleTape_decidesInTime}
  \leanok
  \uses{def:ntm-decides}
  If a nondeterministic machine $N$ with $k \ge 1$ work tapes decides $L$ in
  time $T$, then its single-tape simulator, a machine with one work tape,
  decides $L$ in time $16(k+1)(T(n) + n + 1)^2$. Consequently, for any $k$, if
  $L$ is decided by some $k$-work-tape machine in time $T = O(n^c)$, it is
  decided by some one-work-tape machine in time $O(n^{c'})$ for some $c'$.
\end{theorem}
\begin{proof}
  \leanok
  Interleave the $k$ work tapes on one tape in blocks carrying head markers;
  each simulated step sweeps the used region. For $k = 0$, pad with a dummy
  work tape.
\end{proof}

\begin{theorem}[Single-tape reduction, deterministic]
  \label{thm:models-single-tape}
  \lean{Complexity.TM.exists_singleTape_decidesInTime}
  \leanok
  \uses{def:tm-decides}
  If a deterministic machine with $k$ work tapes decides $L$ in time $T$, then
  some deterministic machine with one work tape decides $L$ in time
  $16(k+1)(T(n) + n + 1)^2$.
\end{theorem}
\begin{proof}
  \leanok
  \uses{lem:models-to-ntm, thm:models-ntm-single-tape}
  Embed as a nondeterministic machine, apply the single-tape simulation, which
  preserves determinism and the output-$0$ rejection convention, and convert
  back.
\end{proof}

\section{Universal simulation}

\begin{definition}[Simulation]
  \label{def:simulates}
  \lean{Complexity.TM.Simulates}
  \leanok
  \uses{def:produces}
  A deterministic machine $U$ simulates $M$ under a compiler
  $c : \{0,1\}^* \to \{0,1\}^*$ if for every program $p$: $U$ halts on $c(p)$
  if and only if $M$ halts on $p$, and for every $y$, $U$ produces $y$ on $c(p)$
  if and only if $M$ produces $y$ on $p$.
\end{definition}

\begin{definition}[Timed simulation and overhead policies]
  \label{def:models-simulates-in-time}
  \lean{Complexity.TM.SimulatesInTime, Complexity.TM.HasAdditiveProgramOverhead, Complexity.TM.PolynomialTimeOverhead}
  \leanok
  \uses{def:produces}
  $U$ simulates $M$ under $c$ in time $\tau : \{0,1\}^* \times \mathbb{N} \to \mathbb{N}$
  if whenever $M$ halts on $p$ (resp.\ produces $y$ on $p$) within $t$ steps,
  $U$ halts on $c(p)$ (resp.\ produces $y$ on $c(p)$) within $\tau(p, t)$ steps.
  The compiler has additive overhead $C$ if $|c(p)| \le |p| + C$ for all $p$.
  The clock $\tau$ is polynomial if $\tau(p, t) \le a (|p| + t + 1)^e$ for
  some constants $a, e$.
\end{definition}

\begin{definition}[Computable function]
  \label{def:computable}
  \lean{Complexity.TM.IsComputable}
  \leanok
  \uses{def:tm-computes}
  A function $f : \{0,1\}^* \to \{0,1\}^*$ is computable if some deterministic
  machine, with any number of work tapes, computes $f$.
\end{definition}

\begin{definition}[Universal machine]
  \label{def:universal}
  \lean{Complexity.TM.IsUniversal, Complexity.TM.IsEfficientlyUniversal, Complexity.TM.IsEfficientlyUniversalFor}
  \leanok
  \uses{def:simulates, def:models-simulates-in-time, def:computable}
  A deterministic machine $U$ is universal if for every $k$ and every machine
  $M$ with $k$ work tapes there is a computable compiler $c$ such that $U$
  simulates $M$ under $c$. It is efficiently universal for an admissibility
  predicate on clocks if for every such $M$ there are a computable compiler
  $c$, a constant $C$, and a clock $\tau$ such that $U$ simulates $M$ under
  $c$, $c$ has additive overhead $C$, $U$ simulates $M$ under $c$ in time
  $\tau$, and $\tau$ is admissible; it is efficiently universal when the
  predicate is polynomiality. Computability of $c$ is essential: without it a
  machine that copies its input after a leading $0$ and diverges after a
  leading $1$ would be universal, with a compiler that consults the source
  machine's halting behavior and output.
\end{definition}

\begin{lemma}[Simulation algebra]
  \label{lem:models-simulation-algebra}
  \lean{Complexity.TM.Simulates.refl, Complexity.TM.Simulates.comp, Complexity.TM.SimulatesInTime.comp, Complexity.TM.HasAdditiveProgramOverhead.comp, Complexity.TM.PolynomialTimeOverhead.comp, Complexity.TM.IsComputable.comp}
  \leanok
  \uses{def:simulates, def:models-simulates-in-time, def:computable}
  Every machine simulates itself under the identity. If $U_1$ simulates $U_2$
  under $c_1$ and $U_2$ simulates $M$ under $c_2$, then $U_1$ simulates $M$
  under $c_1 \circ c_2$, and likewise for timed simulations with clock
  $(p, t) \mapsto \tau_1(c_2(p), \tau_2(p, t))$. Additive overheads add;
  polynomial clocks compose through a compiler with additive overhead; and
  computable functions are closed under composition.
\end{lemma}
\begin{proof}
  \leanok
  Unfold the definitions; composition of computable functions sequences two
  machines.
\end{proof}

\begin{lemma}[Universality transfers]
  \label{lem:models-universal-transfer}
  \lean{Complexity.TM.Simulates.isUniversal, Complexity.TM.IsEfficientlyUniversal.isUniversal, Complexity.TM.IsEfficientlyUniversalFor.isUniversal}
  \leanok
  \uses{def:universal, lem:models-simulation-algebra}
  A machine that simulates a universal machine under a computable compiler is
  universal. Every machine that is efficiently universal for some
  admissibility predicate, in particular every efficiently universal machine,
  is universal.
\end{lemma}
\begin{proof}
  \leanok
  \uses{lem:models-simulation-algebra}
  Compose compilers for the first claim; for the second, forget the overhead
  data.
\end{proof}

\begin{definition}[The universal machine]
  \label{def:models-utm}
  \lean{Complexity.TM.utmTM}
  \leanok
  \uses{def:models-combinators}
  The universal machine $\mathcal{U}$ has six work tapes and is
  $\mathrm{seq}(\mathrm{init}, \mathrm{seq}(\mathrm{loop}(\mathrm{body},
  \mathrm{haltTest}), \mathrm{extract}))$. On input $\langle \alpha, x \rangle$
  it parses the description $\alpha$ of a one-work-tape machine, simulates one
  step per loop iteration, tests for the halt state, and finally copies the
  simulated output tape to its own output tape.
\end{definition}

\begin{theorem}[Universal simulation of deciders]
  \label{thm:models-utm-deciders}
  \lean{Complexity.TM.UTMBody.utmTM_universal, Complexity.TM.UTMBody.utmTM_universal_padded, Complexity.TM.UTMBody.utmTM_simulates_decider}
  \leanok
  \uses{def:models-utm, def:tm-decides, def:pair}
  Let $M$ be a deterministic machine with $k$ work tapes that decides $L$ in
  time $T$, and let $T'(n) = 16(k+1)(T(n) + n + 1)^2$. There is a description
  $\alpha_0$ such that for every string $j$ and every input $x$, with
  $\alpha = \alpha_0 j$, the machine $\mathcal{U}$ on
  $\langle \alpha, x \rangle$ halts within
  $\mathrm{utmTime}(\alpha, T'(|x|), |x|)$ steps with output cell $1$ equal to
  $1$ if $x \in L$ and $0$ if $x \notin L$. Here
  $\mathrm{utmTime}(\alpha, t, n) = 4(2|\alpha| + 2 + n) + 4 g(\alpha) +
  (t+1) s(\alpha) + 2t + 35$, where $g(\alpha)$ and $s(\alpha)$ depend only on
  $\alpha$ (Arora--Barak Theorem~1.9, with the quadratic factor coming from the
  single-tape reduction). The unpadded case, $j$ empty, is also stated on its
  own. The underlying linear-time statement is also formalized: for every
  description $\alpha$ satisfying a syntactic side condition on its table
  region (met by every encoded machine, with or without padding), if the
  one-work-tape machine that $\alpha$ describes decides $L$ in time $T$, then
  $\mathcal{U}$ on $\langle \alpha, x \rangle$ halts within
  $\mathrm{utmTime}(\alpha, T(|x|), |x|)$ steps with the same verdict.
\end{theorem}
\begin{proof}
  \leanok
  \uses{thm:models-single-tape, thm:models-seq-hoare, thm:models-loop-hoare}
  Reduce $M$ to one work tape, encode it, and compose Hoare triples for the
  parsing, per-step simulation loop, and extraction phases. Descriptions
  tolerate arbitrary trailing padding.
\end{proof}

\begin{theorem}[The universal machine is universal]
  \label{thm:utm}
  \lean{Complexity.TM.utmTM_simulates_pair, Complexity.TM.utmTM_isEfficientlyUniversal, Complexity.TM.utmTM_isUniversal}
  \leanok
  \uses{def:models-utm, def:simulates, def:models-simulates-in-time, def:universal, def:pair}
  For every $k$ and every machine $M$ with $k$ work tapes there is a description
  $\alpha$ such that $\mathcal{U}$ simulates $M$ under the compiler
  $p \mapsto \langle \alpha, p \rangle$, and does so in time
  $(p, t) \mapsto \mathrm{utmTime}(\alpha, 16(k+1)(t + |p| + 1)^2, |p|)$.
  Consequently $\mathcal{U}$ is efficiently universal (with additive program
  overhead $2|\alpha| + 2$) and universal.
\end{theorem}
\begin{proof}
  \leanok
  \uses{thm:models-utm-deciders, lem:models-simulation-algebra, lem:models-universal-transfer, lem:encodings-pair-length, thm:encodings-pair-fp}
  The phase Hoare triples behind Theorem~\ref{thm:models-utm-deciders} are
  applied to arbitrary programs instead of deciders: the single-tape
  simulation of $M$, the interpretation of its encoded description, and
  $\mathcal{U}$ each preserve halting and exact output within their clocks and
  diverge whenever the simulated machine diverges. Composing the three
  simulations gives the clock above. The compiler
  $p \mapsto \langle \alpha, p \rangle$ is polynomial-time computable and adds
  exactly $2|\alpha| + 2$ bits, and the composed clock is bounded by a
  polynomial in $|p| + t + 1$.
\end{proof}

\begin{theorem}[Undecidability of the halting problem]
  \label{thm:models-halting-undecidable}
  \uses{def:produces, def:tm-decides, def:universal, def:models-utm}
  Let $H = \{ p : \mathcal{U} \text{ halts on } p \}$. There are no $k$,
  machine $M$ with $k$ work tapes, and bound $T$ such that $M$ decides $H$ in
  time $T$. The same holds for $H_V$ for every universal machine $V$.
\end{theorem}
\begin{proof}
  \uses{thm:utm, lem:models-universal-transfer}
  Diagonalize: from a decider for $H$, build a machine that halts on its own
  description exactly when $\mathcal{U}$ does not, and compile it through the
  universality compiler.
\end{proof}

\section{The time hierarchy}

\begin{definition}[Time-constructible function]
  \label{def:models-time-constructible}
  \lean{Complexity.TimeConstructible}
  \leanok
  \uses{def:tm-computes}
  $T : \mathbb{N} \to \mathbb{N}$ is time-constructible if $T(n) \ge n$ for all
  $n$ and some deterministic machine computes $x \mapsto \mathrm{bin}(T(|x|))$
  (least significant bit first) in time $O(T(n))$ (Arora--Barak
  Definition~1.12). The formalized hierarchy theorem does not use this notion.
\end{definition}

\begin{definition}[Clock-constructible function]
  \label{def:models-clock-constructible}
  \lean{Complexity.TM.ClockConstructible}
  \leanok
  \uses{def:models-hoare}
  $g$ is clock-constructible if there are an eight-work-tape machine $M$ and a
  constant $C$ such that, for every input $x$ and every starting layout in
  which the input tape holds $x$ with its head at cell $1$, every work head is
  at a cell $\ge 1$ not reading $\triangleright$, work tapes $5$ and $6$ are
  blank with their heads at cell $1$, and the output head is at cell $1$ with
  $\triangleright$ at cell $0$ and nowhere else, $M$ started in its start state
  halts within $C (g(|x|) + |x| + 1)$ steps with: work tape $6$ holding the
  unary clock of length $g(|x|)$ (cells $1, \dots, g(|x|)$ equal to $1$, the
  rest blank, head at cell $1$); every other work tape exactly as it started
  (so the scratch tape $5$ is again blank); the input head back at cell $1$;
  and the output head at cell $1$ with $\triangleright$ still at cell $0$ and
  nowhere else. The layout is pinned to the diagonalizer that consumes it.
\end{definition}

\begin{lemma}[Polynomial clocks]
  \label{lem:models-clock-pow}
  \lean{Complexity.TM.clockConstructible_succ, Complexity.TM.ClockConstructible.mul_succ, Complexity.TM.clockConstructible_pow}
  \leanok
  \uses{def:models-clock-constructible}
  The function $n \mapsto n + 1$ is clock-constructible, and if $g$ is
  clock-constructible then so is $n \mapsto g(n)(n+1)$. Hence for every
  $k \ge 1$ the function $n \mapsto (n+1)^k$ is clock-constructible.
\end{lemma}
\begin{proof}
  \leanok
  $n + 1$ is constructed by one input sweep. For the product, run the clock
  for $g$, move it to the scratch tape, and append it to the clock tape once
  per input position. Induction on $k$ gives the powers.
\end{proof}

\begin{theorem}[Deterministic time hierarchy, weak form]
  \label{thm:time-hierarchy}
  \lean{Complexity.time_hierarchy_weak, Complexity.time_hierarchy_weak_ssubset, Complexity.DTIME_pow_ssubset}
  \leanok
  \uses{def:dtime, def:models-clock-constructible}
  Let $g$ be clock-constructible with $g(n) \ge 1$ for all $n$, and let
  $(f(n) + n + 1)^2 = o(g(n))$. Then some language lies in
  $\cc{DTIME}((n+1)^2 (g(n) + 1))$ but not in $\cc{DTIME}(f)$; in fact
  $\cc{DTIME}(f) \subsetneq \cc{DTIME}((n+1)^2 (g(n) + 1))$. In particular, for
  every $a \ge 1$,
  $\cc{DTIME}((n+1)^a) \subsetneq \cc{DTIME}((n+1)^{2a+5})$.
\end{theorem}
\begin{proof}
  \leanok
  \uses{def:models-utm, thm:models-seq-hoare, thm:models-if-hoare, thm:models-single-tape, lem:models-clock-pow}
  The diagonalizer is an eight-work-tape machine built with the combinators.
  On input $x$ it forms $\langle x, x \rangle$ on a work tape, outputs $0$ if
  $x$ fails the syntactic side condition on descriptions, and otherwise runs
  the clock for $g(|x|)$, runs a clocked variant of $\mathcal{U}$ (built from
  the same initialization, halt-test, and extraction phases) on
  $\langle x, x \rangle$ within that budget, and negates the verdict; it runs
  within a constant times $(n+1)^2 (g(n) + 1)$ steps. A hypothetical decider in
  time $f_0 = O(f)$ is reduced to one work tape, and its description, padded
  with enough zeros, is an input on which the simulation budget
  $16(k+1)(f_0(n) + n + 1)^2$ falls below $g(n)$; there the diagonalizer's
  verdict contradicts the decider's. The polynomial case takes
  $g(n) = (n+1)^{2a+3}$.
\end{proof}

\begin{theorem}[Deterministic time hierarchy, textbook form]
  \label{thm:models-time-hierarchy-strong}
  \uses{def:dtime, def:models-time-constructible}
  If $f$ and $g$ are time-constructible and $f(n) \log f(n) = o(g(n))$, then
  $\cc{DTIME}(f) \subsetneq \cc{DTIME}(g)$ (Arora--Barak Theorem~3.1).
\end{theorem}
\begin{proof}
  \uses{thm:utm}
  Needs a universal simulation with $O(T \log T)$ overhead
  (Hennie--Stearns), which the current machine does not provide: its overhead
  is quadratic through the single-tape reduction. It also needs a bridge from
  time-constructibility to the pinned clock layout.
\end{proof}

\section{Random access machines}

The random access machine is a register machine with indirect addressing,
measured by \emph{logarithmic} cost. Its role is to show that $\cc{P}$ does not
depend on the Turing-machine conventions.

\begin{definition}[Random access machine]
  \label{def:models-ram}
  \lean{Complexity.RAM.Instr, Complexity.RAM.Program, Complexity.RAM.Cfg, Complexity.RAM.step, Complexity.RAM.run, Complexity.RAM.initCfg, Complexity.RAM.Instr.logCost, Complexity.RAM.logTimeUpto, Complexity.RAM.unitTimeUpto, Complexity.RAM.spaceUpto}
  \leanok
  A program is a list of instructions over registers $R_0, R_1, \dots$ holding
  natural numbers: immediate load, $\mathrm{add}$, truncated $\mathrm{sub}$,
  $\mathrm{mul}$, indirect $\mathrm{load}$ ($R_d := R_{R_a}$) and
  $\mathrm{store}$ ($R_{R_a} := R_s$), conditional and unconditional jumps, and
  $\mathrm{halt}$; a program counter past the end halts. On input $x$,
  $R_0 = |x|$, $R_{i+1} = x_i$, and all other registers are $0$. Each
  instruction costs $1$ plus bit lengths of the values it touches: the
  constant of an immediate load; both operands and the result for $\mathrm{add}$
  and $\mathrm{mul}$; both operands for $\mathrm{sub}$ (whose result is no
  longer than its first operand); the runtime address and the value moved for
  indirect $\mathrm{load}$ and $\mathrm{store}$; the tested register for a
  conditional jump; nothing more for $\mathrm{jmp}$ and $\mathrm{halt}$.
  Register indices and jump targets written in the program are not charged.
  The logarithmic time of a run is the total cost of its non-halted steps, and
  its unit time is the number of those steps. Space is the peak, over the
  configurations visited, of the total bit length of the addresses and
  contents of the nonzero registers.
\end{definition}

\begin{definition}[RAM deciders and classes]
  \label{def:models-ram-classes}
  \lean{Complexity.RAM.Program.DecidesInTime, Complexity.RAM.Program.DecidesInSpace, Complexity.RAM.DTIME, Complexity.RAM.DSPACE, Complexity.RAM.P}
  \leanok
  \uses{def:models-ram, def:language}
  A program decides $L$ in time $T$ if on every input $x$ it halts after
  logarithmic cost at most $T(|x|)$ with $R_0 = 1$ if $x \in L$ and $R_0 = 0$
  otherwise; deciding in space $S$ bounds peak space instead. Then
  $\cc{DTIME}_{\mathrm{RAM}}(T)$ and $\cc{DSPACE}_{\mathrm{RAM}}(S)$ are the
  languages decided in time $O(T)$ and space $O(S)$, and
  $\cc{P}_{\mathrm{RAM}} = \bigcup_k \cc{DTIME}_{\mathrm{RAM}}(n^k)$.
\end{definition}

\begin{theorem}[Unit cost is unsound]
  \label{thm:models-ram-log-gap}
  \lean{Complexity.RAM.logGap_squaring}
  \leanok
  \uses{def:models-ram}
  For every $k \ge 1$ there are a program and a configuration whose run halts
  after $k + 1$ steps, with unit time exactly $k + 1$ and logarithmic time at
  least $2^k$.
\end{theorem}
\begin{proof}
  \leanok
  Load $2$ into a register and square it $k$ times; the last multiplication
  writes $2^{2^k}$, which has $2^k + 1$ bits.
\end{proof}

\begin{theorem}[Turing machines to RAMs]
  \label{thm:models-p-subset-ram-p}
  \lean{Complexity.RAM.TMConfig.Sparse.P_subset_RAM_P, Complexity.RAM.TMConfig.Sparse.compiledDecision_decidesInTime, Complexity.RAM.TMConfig.Sparse.mem_DTIME_of_decidesInTime}
  \leanok
  \uses{def:models-ram-classes, def:P}
  $\cc{P} \subseteq \cc{P}_{\mathrm{RAM}}$. More precisely, if a
  deterministic machine $M$ decides $L$ in time $T$, a fixed RAM program
  built from $M$ decides $L$ in logarithmic time $\beta_M(n, T(n))$, for an
  explicit bound $\beta_M(n, t)$ depending only on $M$, $n$, and $t$; so
  $L \in \cc{DTIME}_{\mathrm{RAM}}(n \mapsto \beta_M(n, T(n)))$.
\end{theorem}
\begin{proof}
  \leanok
  One fixed RAM program, depending only on the Turing machine, stores the
  configuration in a sparse interleaved register layout, computes cell
  addresses at run time, and follows the halting run before copying the verdict
  to $R_0$. When $T$ is bounded by a polynomial, so is
  $n \mapsto \beta_M(n, T(n))$.
\end{proof}

\begin{theorem}[RAMs to Turing machines, parametric]
  \label{thm:models-ram-dtime-sq}
  \lean{Complexity.RAM.RegisterStore.Machine.RAM_DTIME_subset_DTIME_sq, Complexity.RAM.RegisterStore.Machine.denseProgramDecision_decidesInTime}
  \leanok
  \uses{def:models-ram-classes, def:dtime}
  If $n + 1 = O(T(n))$, then
  $\cc{DTIME}_{\mathrm{RAM}}(T) \subseteq \cc{DTIME}(T(n)^2)$. More
  precisely, with no hypothesis on $T$: if a program $P$ decides $L$ in
  logarithmic time $T$, a fixed twenty-work-tape machine built from $P$
  decides $L$ in time $c_P (n + T(n) + 1)^2$, where $c_P$ depends only on $P$.
\end{theorem}
\begin{proof}
  \leanok
  A twenty-work-tape machine built from the program keeps the public input
  immutable, stores the mutable registers as a sparse tagged overlay of
  address/value pairs, and charges each simulated step by the width of its
  instruction; the charges sum to $O((n + T(n))^2)$.
\end{proof}

\begin{theorem}[RAMs to Turing machines, polynomial]
  \label{thm:models-ram-p-subset-p}
  \lean{Complexity.RAM.RegisterStore.Machine.RAM_P_subset_P, Complexity.RAM.RegisterStore.Machine.programDecision_decidesInTime}
  \leanok
  \uses{def:models-ram-classes, def:P}
  $\cc{P}_{\mathrm{RAM}} \subseteq \cc{P}$. The proof goes through a second
  simulator: if a program $P$ decides $L$ in logarithmic time $T$, a fixed
  twenty-work-tape machine built from $P$ decides $L$ in time
  $10^9 m_P (n + (m_P + 2) T(n) + m_P + 4)^4$, where $m_P$ depends only on $P$.
\end{theorem}
\begin{proof}
  \leanok
  For each program, a twenty-work-tape machine built from it simulates the
  program within a fourth-degree polynomial in input length and logarithmic
  time; the number of RAM steps is at most the logarithmic time.
\end{proof}

\begin{theorem}[Polynomial time is model-independent]
  \label{thm:models-ram-p-eq-p}
  \lean{Complexity.RAM.RegisterStore.Machine.RAM_P_eq_P}
  \leanok
  \uses{def:models-ram-classes, def:P}
  $\cc{P}_{\mathrm{RAM}} = \cc{P}$.
\end{theorem}
\begin{proof}
  \leanok
  \uses{thm:models-p-subset-ram-p, thm:models-ram-p-subset-p}
  Antisymmetry of the two inclusions.
\end{proof}

\begin{theorem}[Turing machines to RAMs, parametric]
  \label{thm:models-dtime-subset-ram-dtime}
  \uses{def:models-ram-classes, def:dtime}
  Under an explicit hypothesis that $T$ dominates the input length,
  $\cc{DTIME}(T) \subseteq \cc{DTIME}_{\mathrm{RAM}}(T(n) \log T(n))$.
  What is formalized is only the exact-bound form of
  Theorem~\ref{thm:models-p-subset-ram-p}, membership in
  $\cc{DTIME}_{\mathrm{RAM}}(n \mapsto \beta_M(n, T(n)))$; the asymptotic
  estimate of $\beta_M$ is open.
\end{theorem}
\begin{proof}
  \uses{thm:models-p-subset-ram-p}
  In the explicit bound $\beta_M(n, t)$ of the fixed simulator, each of the
  $t + 1$ simulated steps is charged a constant number of words whose width is
  the bit length of a bound polynomial in $n + t$, so the core should cost
  $O(t \log(n + t))$; the input marshaller's cost must be bounded as well and
  absorbed by the domination hypothesis.
\end{proof}

\begin{theorem}[Polynomial space is model-independent]
  \label{thm:models-ram-pspace}
  \uses{def:models-ram-classes, def:PSPACE}
  $\bigcup_k \cc{DSPACE}_{\mathrm{RAM}}(n^k) = \cc{PSPACE}$.
\end{theorem}
\begin{proof}
  \uses{thm:models-ram-dtime-sq, thm:models-p-subset-ram-p}
  Bound the space of the two existing simulations, charging the sparse
  register store and the simulated tape contents.
\end{proof}

\begin{theorem}[Register-machine and Boolean-program variants]
  \label{thm:models-ram-variants}
  \uses{def:models-ram-classes, def:P}
  \notready
  Define register machines without indirect addressing and Boolean
  (straight-line or bounded-loop) programs, and relate each to the RAM and to
  Turing machines by explicit simulations, so that their polynomial-time
  classes equal $\cc{P}$ as theorems rather than as separate class
  definitions.
\end{theorem}
